\section{Corpus Identity Scheme}
\label{sec:appendix-identity}

\subsection{The URN grammars}

Identity is \code{uuid5} over a canonical \mbox{URN}. One grammar per work:

\begin{center}\small
\begin{tabular}{@{}ll@{}}
\toprule
Work & Mantra \mbox{URN} template \\
\midrule
\d{R}gveda \'S\=akala &
  \code{urn:vedagraph:mantra:rigveda:shakala:} \\
& \quad\code{mandala:\{m\}:sukta:\{s\}:mantra:\{v\}} \\[2pt]
Yajurveda \textsc{vsm} &
  \code{urn:vedagraph:mantra:yajurveda:} \\
& \quad\code{vajasaneyi-madhyandina:adhyaya:\{a\}:mantra:\{v\}} \\[2pt]
Atharvaveda \'Saunaka &
  \code{urn:vedagraph:mantra:atharvaveda:shaunaka:} \\
& \quad\code{kanda:\{k\}:sukta:\{s\}:mantra:\{v\}} \\[2pt]
S\=amaveda Kauthuma &
  \code{urn:vedagraph:mantra:samaveda:kauthuma:\{collection\}} \\
& \quad\code{[:prapathaka:\{p\}][:ardha:\{r\}][:dasati:\{d\}]:verse:\{v\}} \\
\bottomrule
\end{tabular}
\end{center}

Note the S\=amavedic grammar names its collection rather than numbering it
(\cref{sec:corpus}): the first ordinal slot denotes different collections in
different witnesses, and the blast radius of the rejected reading was measured at
1{,}280 of 1{,}875 verses before the decision was taken.

Keys are padded for lexical sorting (\code{VG:RV:SAK:M01:S001:V001}); \mbox{URN}s are
unpadded and are the hash input. Only the \mbox{URN} may never be re-spelled.

\subsection{Verification}

Recomputing every identifier from its \mbox{URN} over all four canonical builds:

\begin{center}\small
\begin{tabular}{@{}lrrrrr@{}}
\toprule
Build & Passages & Dup.\ keys & Dup.\ \mbox{URN}s & Dup.\ \mbox{UUID}s & Mismatches \\
\midrule
\d{R}gveda            & 11{,}590 & 0 & 0 & 0 & 0 \\
S\=amaveda \=Arcika   &  2{,}342 & 0 & 0 & 0 & 0 \\
V\=ajasaneyi          &  2{,}015 & 0 & 0 & 0 & 0 \\
Atharvaveda (working) &  6{,}590 & 0 & 0 & 0 & 0 \\
\midrule
\textbf{pooled}       & \textbf{22{,}537} & \textbf{0} & \textbf{0} & \textbf{0} & \textbf{0} \\
\bottomrule
\end{tabular}
\end{center}

\subsection{The comparison ladder}

Six named comparison forms, from the identity function on stored text to the
cross-script surface: \code{SOURCE\_ORIGINAL}, \code{NFC},
\code{ACCENT\_PRESERVING\_NORMALIZED}, \code{ACCENT\_STRIPPED\_COMPARISON},
\code{SEARCH\_NORMALIZED}, \code{CROSS\_SCRIPT\_COMPARISON}. Stored text is never
replaced: \code{text\_original} and \code{text\_nfc} are separate fields, and
\code{SOURCE\_ORIGINAL} is the identity function, verified by test.

Accent stripping removes only Vedic tone marks (U+0951--U+0954, the combining members
of the Vedic Extensions block, plus U+0300, U+0301, U+030D, U+0331, U+032D), and
adjudicates ambiguous combining marks per base letter. One consequence is worth
stating, because it bears on the flagship case study: the S\=amavedic svara marks
are \emph{not} in those ranges. Every one of the 25{,}542 marks on the 1{,}136
accented witnesses sits in Devanagari Extended (U+A8E1--U+A8F3 --- combining
Devanagari digits one, two and three, plus the combining letters), and zero are in
U+0951--U+0954. The implementation is unaffected, because the S\=amavedic search
surface is reached by transliteration rather than by this fold; the enumeration
above is nonetheless not a description of the accents this corpus actually carries: on a vowel base a tone-set
mark is an accent; on a consonant base carrying exactly one such mark, the mark is
part of the letter (\iast{\'s} is \emph{s} plus U+0301); on a consonant base carrying
more than one, the extra is the tone.

\subsection{Deliberate non-folds and one deliberate loss}

Avagraha is a real orthographic sign and is not folded, at a measured cost of 46 of
116 residual unclassifiable Yajurvedic pairs. A private-use code point occurring 20
times in one witness is not folded, because reading it as an anusv\=ara from context
is a transcription judgement rather than a normalisation; it is registered as a
source defect instead. Conversely the whole lateral series folds onto one sentinel
because U+1E37 denotes the vocalic \iast{l} in one source convention and the
retroflex \iast{\d{l}} in another. Splitting it was implemented, measured, and
reverted.

Case is folded \emph{before} the transcription table, not after: folding the other
way left upper-case variants unmatched and made the search surface non-idempotent.
Reordering was measured to change the surface for 0 of 21{,}936 canonical text
records.

\subsection{The referent-drift guard}

Two digests per key: one over the bytes the source spells, one over a
normalisation-independent comparison surface. Both are compared against the previous
release's digests for the same key, and a difference is adjudicated into
\code{UNCHANGED\_REFERENT}, \code{INTENTIONAL\_REFERENT\_CORRECTION},
\code{REFERENT\_DRIFT} (blocks release), \code{REMOVED\_INVALID\_PASSAGE} or
\code{NEWLY\_DISCOVERED\_PASSAGE}.

Text equality alone is insufficient. Of 1{,}844 released S\=amavedic keys only
1{,}651 have a distinct comparison digest: 192 equivalence classes covering 385 keys
share one, and 174 of those classes share a byte-identical raw digest. The guard
therefore requires three fields to agree --- comparison digest, source locator and
source verse marker --- before calling two observations the same occurrence.

Because the comparison digest runs through the search surface, adding a rule to the
transcription table changes the digest of already-released keys and the guard reports
drift where none occurred. This was measured (one rule moves five Yajurvedic keys,
others four and one more) and resolved by splitting the fold tables: a frozen table
inside the digest, and a free table for comparison only.
